\section*{Capítulo \ref{ch:b2:mantle-melting} --- La fusión del manto}

\begin{solution}{exo:b2:mantle-melting:1}
(a) Más de un \qty{90}{\percent} de olivino: \omterm{def:b2:mantle-melting:peridotites}{dunita}. (b) Olivino y ortopiroxeno
sin clinopiroxeno: \omterm{def:b2:mantle-melting:peridotites}{harzburgita}. (c) Más de un \qty{40}{\percent} de olivino con
los dos piroxenos: \omterm{def:b2:mantle-melting:peridotites}{lherzolita}. (d) Menos de un \qty{40}{\percent} de olivino: una piroxenita
(una websterita olivínica, con los dos piroxenos). (e) $57/98 = \qty{58.2}{\percent}$,
$28/98 = \qty{28.6}{\percent}$, $13/98 = \qty{13.3}{\percent}$: una \omterm{def:b2:mantle-melting:peridotites}{lherzolita}.
\end{solution}

\begin{solution}{exo:b2:mantle-melting:2}
(a) El liquidus con un \qty{60}{\percent} de anortita está a \qty{1390}{\celsius}, en el
campo de la anortita: la anortita cristaliza primero. (b) El diópsido aparece en el
eutéctico, a \qty{1270}{\celsius}, cuando el líquido tiene un \qty{42}{\percent} de anortita.
(c) El primer líquido de cualquier mezcla de los dos minerales es el líquido eutéctico,
con un \qty{42}{\percent} de anortita, a \qty{1270}{\celsius}.
\end{solution}

\begin{solution}{exo:b2:mantle-melting:3}
(a) $\num{4.0e-5}\times1623/(3300\times1000) = \qty{1.97e-8}{K/Pa} =
\qty{19.7}{K/GPa}$. (b) $\num{e9}/(3300\times9.81) = \qty{3.09e4}{m}$, o sea
\qty{30.9}{km} por gigapascal; $19.7/30.9 = \qty{0.64}{K/km}$, unas cincuenta veces
menos que los \qty{30}{K/km} cerca de la superficie, donde el calor se conduce y no se
transporta.
\end{solution}

\begin{solution}{exo:b2:mantle-melting:4}
(a) $T_{\mathrm s}(0) = \qty{1085.7}{\celsius}$; $T_{\mathrm s}(1) = 1085.7 +
132.9 - 5.1 = \qty{1213.5}{\celsius}$; $T_{\mathrm s}(2) = 1085.7 + 265.8 - 20.4 =
\qty{1331.1}{\celsius}$; $T_{\mathrm s}(3) = 1085.7 + 398.7 - 45.9 =
\qty{1438.5}{\celsius}$. (b) $132.9 - 2\times5.1\times2 = \qty{112.5}{K/GPa}$.
(c) A \qty{3}{GPa}, \qty{1400}{\celsius} está por debajo del solidus,
\qty{1438.5}{\celsius}; al subir hasta \qty{2}{GPa} por la adiabata se enfría solo
unos \qty{20}{\celsius}, hasta cerca de \qty{1380}{\celsius}, mientras el solidus baja a
\qty{1331.1}{\celsius}: la roca queda por encima de su solidus y una parte se funde.
\end{solution}

\begin{solution}{exo:b2:mantle-melting:5}
(a) $1350 + (1277.5 - 1350)\times(30 - 20)/(40 - 20) = \qty{1314}{\celsius}$.
(b) La composición está del lado del diópsido respecto al eutéctico: diópsido.
(c) El líquido alcanza el eutéctico con un \qty{42}{\percent} de anortita; la anortita
se queda en el líquido, así que la fracción líquida es $30/42 = 0.714$ y un $1 - 0.714 =
\qty{29}{\percent}$ del sistema ha cristalizado como diópsido. (d) Todo
cristaliza: \qty{70}{\percent} de diópsido y \qty{30}{\percent} de anortita, la
composición global.
\end{solution}

\begin{solution}{exo:b2:mantle-melting:6}
(a) Forsterita: $60.08/(2\times40.30 + 60.08) = 60.08/140.68 =
\qty{42.71}{\percent}$ de \ce{SiO2}; enstatita: $60.08/100.38 = \qty{59.85}{\percent}$.
(b) $0.975\times59.85 + 2.5 = \qty{60.85}{\percent}$. (c) La enstatita está entre la
forsterita y el líquido: la fracción de forsterita es $(60.85 - 59.85)/(60.85 -
42.71) = 1.00/18.14 = \qty{5.5}{\percent}$, el valor de Bowen y Andersen. (d) El
fundido es más rico en sílice que el ortopiroxeno del que procede, y el residuo
gana olivino.
\end{solution}

\begin{solution}{exo:b2:mantle-melting:7}
(a) $5.1P^2 - 112.9P + 214.3 = 0$: $P_0 = (112.9 - \sqrt{12746.4 -
4371.7})/10.2 = (112.9 - 91.51)/10.2 = \qty{2.10}{GPa}$, a $2.097\times30.9 =
\qty{64.8}{km}$. (b) $5.1P^2 - 112.9P + 364.3 = 0$: $P_0 = (112.9 -
72.90)/10.2 = \qty{3.92}{GPa}$, a \qty{121}{km}. (c) $T_{\mathrm s}'(2.10) = 132.9
- 21.4 = \qty{111.5}{K/GPa}$, luego $\Delta P_0 = 10/(111.5 - 20) = \qty{0.109}{GPa}$
por cada \qty{10}{\celsius}, o sea \qty{3.4}{km}.
\end{solution}

\begin{solution}{exo:b2:mantle-melting:8}
(a) $T_{\mathrm s} = \qty{1331.1}{\celsius}$, $T_{\mathrm l} = 1475.0 + 160.0 -
12.8 = \qty{1622.2}{\celsius}$; $T = 1331.1 + 0.10\times291.1 =
\qty{1360.2}{\celsius}$, o sea \qty{1633}{K}. (b) $(\partial T/\partial P)_F = 112.5
+ 0.10\times(67.2 - 112.5) = \qty{108.0}{K/GPa}$; $(\partial T/\partial F)_P =
T_{\mathrm l} - T_{\mathrm s} = \qty{291.1}{K}$. (c) $c_p/T = 0.612$; con $\alpha/\rho
= \qty{12.1}{J/(kg.K)}$ por gigapascal, el numerador es $0.612\times108.0 - 12.1
= 54.0$ y el denominador $300 + 0.612\times291.1 = 478.2$: $-\mathrm dF/\mathrm
dP = 0.113$, un \qty{11.3}{\percent} por gigapascal. (d) $54.0/178.2 = 0.303$,
un \qty{30}{\percent} por gigapascal: el calor latente divide la productividad por un
factor de 2.7.
\end{solution}

\begin{solution}{exo:b2:mantle-melting:9}
(a) $43\times1 = \qty{43}{\celsius}$; $43\times5^{0.75} = 43\times3.34 =
\qty{144}{\celsius}$; $43\times13^{0.75} = 43\times6.85 = \qty{294}{\celsius}$.
(b) $1213.5 - 294 = \qty{919}{\celsius}$. (c) \qty{125}{ppm} son un
\qty{0.0125}{\percent}, así que $X = 0.0125/0.01 = \qty{1.25}{\percent}$ y $\Delta T =
43\times1.25^{0.75} = 43\times1.18 = \qty{51}{\celsius}$.
\end{solution}

\begin{solution}{exo:b2:mantle-melting:10}
(a) La base mide $2H$ de ancho y el manto la atraviesa a $U$: $2UH$ por unidad de longitud
de dorsal y de tiempo. Cada lado abarca la altura $H$ y se atraviesa horizontalmente a
$U$: $UH$ cada uno, $2UH$ en total. Con lados a $\theta$ la base mediría
$2H/\tan\theta$ y la entrada sería $2UH/\tan\theta$, igual a la salida solo si
$\tan\theta = 1$: en este modelo el ángulo lo fija la conservación de la masa. (b) Con
productividad constante la \omterm{def:b2:mantle-melting:fraction}{fracción de fundido} en el vértice es $F_{\max} = \Pi P_0 = \Pi\rho
g z_0$, y $H = z_0$: $h_c = F_{\max}z_0/2 = \Pi\rho g z_0^2/2$. (c) Cada año se forman
\qty{50}{mm} de placa nueva, $\qty{0.05}{m}\times\qty{1000}{m}\times\qty{7100}{m} =
\qty{3.6e5}{m^3}$ de corteza por kilómetro de dorsal; con una \omterm{def:b2:mantle-melting:fraction}{fracción de fundido} media del
\qty{12}{\percent} procede de $\num{3.6e5}/0.12 = \qty{3.0e6}{m^3}$ de manto.
\end{solution}

\begin{solution}{exo:b2:mantle-melting:11}
(a) Cuando una masa $t$ de A ha salido de una unidad de masa de líquido, el líquido contiene $0.6 -
t$ de A, 0.2 de B y 0.2 de C en una masa $1 - t$: $\big((0.6 - t), 0.2,
0.2\big)/(1 - t)$, y B/C sigue valiendo 1. (b) Sobre la cotéctica, $(0.40 - 0.15s, 0.60 -
0.30s, 0.45s)$; B = C da $0.60 - 0.30s = 0.45s$, $s = 0.8$: $L_1$ tiene A 28, B
36, C 36. Entonces $0.2/(1 - t) = 0.36$, $t = 0.44$: un \qty{44}{\percent} del sistema ha
cristalizado como A. (c) Todo el C está en el líquido: $f = 20/45 = 0.44$, un
\qty{44}{\percent} sigue siendo líquido. (d) A: $0.60 - 0.444\times0.25 = 0.49$; B: $0.20 -
0.444\times0.30 = 0.07$; con el líquido, $49 + 7 + 44 = \qty{100}{\percent}$.
\end{solution}

\begin{solution}{exo:b2:mantle-melting:12}
(a) $X_{\mathrm{sat}}(1) = 12 + 1 = \qty{13}{\percent}$, $\Delta T =
\qty{294}{\celsius}$, solidus saturado $1213.5 - 294 = \qty{919}{\celsius}$;
$X_{\mathrm{sat}}(2) = 12\times1.516 + 2 = \qty{20.2}{\percent}$, $\Delta T =
43\times20.2^{0.75} = 43\times9.53 = \qty{410}{\celsius}$, solidus saturado $1331.1
- 410 = \qty{921}{\celsius}$. (b) El fundido disuelve más agua a mayor presión,
y el mayor descenso compensa la subida de \qty{118}{\celsius} del solidus seco.
(c) El solidus debe bajar $1331.1 - 1100 = \qty{231}{\celsius}$: $X =
(231/43)^{1/0.75} = 5.37^{1.333} = \qty{9.4}{\percent}$ de agua, menos que el
\qty{20.2}{\percent} de saturación, así que la cuña puede fundirse.
\end{solution}

\begin{solution}{pb:b2:mantle-melting:1}
\textbf{1.} $1350 + 273 = \qty{1623}{K}$.
\textbf{2.} $\num{4.0e-5}\times1623/(3300\times1000) = \qty{1.97e-8}{K/Pa} =
\qty{19.7}{K/GPa}$, redondeado a \qty{20}{K/GPa}.
\textbf{3.} $\num{e9}/(3300\times9.81) = \qty{30.9}{km}$; $20/30.9 =
\qty{0.65}{K/km}$.
\textbf{4.} $1350 + 20P = 1085.7 + 132.9P - 5.1P^2$, es decir $5.1P^2 - 112.9P +
264.3 = 0$: $P_0 = (112.9 - \sqrt{12746.4 - 5391.7})/10.2 = (112.9 - 85.76)/10.2 =
\qty{2.66}{GPa}$.
\textbf{5.} $2.66\times30.9 = \qty{82}{km}$; $T = 1350 + 20\times2.66 =
\qty{1403}{\celsius}$.
\textbf{6.} $1085.7 + 132.9\times2.66 - 5.1\times2.66^2 = 1085.7 + 353.5 - 36.1 =
\qty{1403}{\celsius}$.
\textbf{7.} $T_{\mathrm s} = 1085.7 + 199.4 - 11.5 = \qty{1273.6}{\celsius}$;
$T_{\mathrm l} = 1475.0 + 120.0 - 7.2 = \qty{1587.8}{\celsius}$; diferencia
\qty{314.2}{K}.
\textbf{8.} $1273.6 + 0.10\times314.2 = \qty{1305.0}{\celsius}$, \qty{1578}{K}.
\textbf{9.} $T_{\mathrm s}' = 132.9 - 15.3 = \qty{117.6}{K/GPa}$, $T_{\mathrm l}' =
80.0 - 9.6 = \qty{70.4}{K/GPa}$: $(\partial T/\partial P)_F = 117.6 +
0.10\times(70.4 - 117.6) = \qty{112.9}{K/GPa}$; $(\partial T/\partial F)_P =
\qty{314.2}{K}$.
\textbf{10.} $c_p/T = 1000/1578 = 0.634$; numerador $0.634\times112.9 - 12.1 = 71.6
- 12.1 = 59.5$; denominador $300 + 0.634\times314.2 = 499.2$: $0.119$, o sea un
\qty{11.9}{\percent} por gigapascal.
\textbf{11.} Con la función ajustada $F = T'^{1.5}$, donde $T'$ es la fracción del camino del solidus al liquidus, la temperatura debe subir mucho con $F$ para formar
el primer uno por ciento de fundido, de modo que $(\partial T/\partial F)_P$ es grande cerca del solidus
y la fusión empieza despacio. $F_{\max} = 0.09\times2.66 = 0.24$, un \qty{24}{\percent}.
\textbf{12.} $0.13/0.50 = 0.26$: el clinopiroxeno se agota con un \qty{26}{\percent} de
fusión. Con un \qty{24}{\percent} el residuo bajo el eje aún conserva un poco:
una \omterm{def:b2:mantle-melting:peridotites}{lherzolita} pobre en clinopiroxeno, al borde de la \omterm{def:b2:mantle-melting:peridotites}{harzburgita}.
\textbf{13.} $h_c = 0.24\times82/2 = \qty{9.8}{km}$.
\textbf{14.} $F_{\max}/2 = \qty{12}{\percent}$.
\textbf{15.} \qty{9.8}{km} es un \qty{38}{\percent} más que \qty{7.1}{km}, por encima del
intervalo de la corteza normal.
\textbf{16.} Una productividad constante da demasiado fundido al manto más profundo; el
modelo ajustado forma despacio el primer fundido, $F(z)$ se mantiene más bajo cerca del inicio, y
$\int F\,\mathrm dz$ es menor.
\textbf{17.} $F_{\max} = \Pi P_0$ y $H = 30.9P_0$ (km): $h_c = \Pi P_0\times30.9P_0/2
= kP_0^2$ con $k = 0.09\times30.9/2 = \qty{1.39}{km/GPa^2}$.
\textbf{18.} $P_0 = \sqrt{7.1/1.39} = \sqrt{5.11} = \qty{2.26}{GPa}$.
\textbf{19.} $2.26\times30.9 = \qty{70}{km}$.
\textbf{20.} $T_{\mathrm s}(2.26) = 1085.7 + 300.4 - 26.0 = \qty{1360.0}{\celsius}$,
y $T_p = 1360.0 - 20\times2.26 = 1360.0 - 45.2 = \qty{1315}{\celsius}$.
\textbf{21.} Para \qty{6.3}{km}: $P_0 = \sqrt{4.53} = \qty{2.13}{GPa}$, $T_{\mathrm s}
= \qty{1345.6}{\celsius}$, $T_p = \qty{1303}{\celsius}$. Para \qty{7.9}{km}: $P_0 =
\sqrt{5.68} = \qty{2.38}{GPa}$, $T_{\mathrm s} = \qty{1373.1}{\celsius}$, $T_p =
\qty{1326}{\celsius}$. La dispersión de la corteza permite $1315 \pm
\qty{12}{\celsius}$.
\textbf{22.} $P_0 = \sqrt{20/1.39} = \sqrt{14.39} = \qty{3.79}{GPa}$; $T_{\mathrm s} =
1085.7 + 503.7 - 73.3 = \qty{1516.1}{\celsius}$; $T_p = 1516.1 - 75.8 =
\qty{1440}{\celsius}$, un exceso de $1440 - 1315 = \qty{125}{\celsius}$.
\textbf{23.} Ambos usan indicios distintos. Las lavas registran el manto más caliente,
en el centro de la pluma; la corteza promedia los fundidos de todo el manto que
atraviesa la región de fusión, en gran parte más frío. El modelo simple también mantiene
constante la productividad, aunque el clinopiroxeno se agota con un \qty{26}{\percent}
mientras que aquí $\Pi P_0 = 0.09\times3.79 = 0.34$: el manto caliente real se funde menos, y necesita
una temperatura más alta para formar la misma corteza.
\textbf{24.} $T_{\mathrm s}'(2.26) = 132.9 - 23.1 = \qty{109.8}{K/GPa}$; $\Delta P_0 =
10/(109.8 - 20) = \qty{0.111}{GPa}$; como $h_c \propto P_0^2$, $\Delta h_c/h_c =
2\times0.111/2.26 = \qty{9.8}{\percent}$, cerca de una décima parte por cada \qty{10}{\celsius}.
\textbf{25.} En una dorsal muy lenta el manto sube tan despacio que pierde calor por
conducción hacia la superficie fría; abandona la adiabata, su fusión se detiene
a más profundidad, y se recorta la parte alta del triángulo, donde $F$ es mayor: la corteza es
más delgada.
\textbf{26.} La corteza de las dorsales normales da una \omterm{def:b2:mantle-melting:tp}{temperatura potencial del manto} de
\textbf{\qty[text-series-to-math=true]{1315}{\celsius}}, con un margen de unos
\qty{12}{\celsius}: \qty{35}{\celsius} por debajo de los \qty{1350}{\celsius} de la física
de minerales, cuya incertidumbre es más o menos la misma, y muy dentro del intervalo de 1280 a
\qty{1400}{\celsius} de las lavas. Tres termómetros independientes coinciden en un manto
de cerca de \qty{1300}{\celsius} a \qty{1350}{\celsius}.
\end{solution}
